\documentclass[11pt]{article}
\usepackage{amsgen,amsmath,amstext,amsbsy,amsopn,amsfonts,amsthm,amscd}
%\usepackage{bbm}
\usepackage[dvips]{graphicx}
\usepackage{enumerate}
\setlength{\textwidth}{16.0cm} \setlength{\textheight}{25cm}
\voffset=-2.50cm \hoffset=-2.00cm \setlength{\topmargin}{0.4cm}

%%%%%%%%%%%%%%%%%%%%%%%%%%%%%%%%%%%%%%%%%%%%%%%%%%%%%%%%%%%%%%

\newtheorem{lem1}{Lema}%[chapter]
\newtheorem{th1}{Teorema}%[chapter]
\newtheorem{prop1}{Proposici\'{o}n}%[chapter]
\newtheorem{defi1}{Definici\'{o}n}%[chapter]
\newtheorem{coro1}{Corolario}%[chapter]
\newtheorem{ob1}{Observaci\'{o}n}%[chapter]
\newtheorem{obs1}{Observaci\'{o}n}%[chapter]
\newtheorem{nota1}{Nota}%[chapter]
\newtheorem{eje1}{Ejemplo}%[chapter]
\newtheorem{ejes1}{Ejemplo}%[chapter]
\newtheorem{ejerc1}{Ejercicio}%[chapter]
\newtheorem{ejercs1}{Ejercicio}%[chapter]
\newenvironment{prop}{\begin{prop1}\rm}{\end{prop1}}
\newenvironment{defi}{\begin{defi1}}{\end{defi1}}
\newenvironment{coro}{\begin{coro1}}{\end{coro1}}
\newenvironment{th2}{\begin{th1}}{\end{th1}}
\newenvironment{nota}{\begin{nota1}\rm}{\end{nota1}}
\newenvironment{lem}{\begin{lem1}}{\end{lem1}}
\newenvironment{ob}{\begin{ob1}\rm}{\end{ob1}}
\newenvironment{obs}{\begin{obs1}\rm}{\end{obs1}}
\newenvironment{eje}{\begin{eje1}\rm}{\end{eje1}}
\newenvironment{ejes}{\begin{ejes1}\rm}{\end{ejes1}}
\newenvironment{ejerc}{\begin{ejerc1}\rm}{\end{ejerc1}}
\newenvironment{ejercs}{\begin{ejercs1}\rm}{\end{ejercs1}}
%%%%%%%%%%%%%%%%%%%%%%%%%%%%%%%%%%%%%%%%%%%%%%%%%%%%%%
%\DeclareSymbolFont{cucu}{U}{bbmss}{m}{n}
\DeclareMathOperator{\CC}{\mathbb C}%{\mathalpha}{cucu}{67}
\DeclareMathOperator{\NN}{\mathbb N}%{\mathalpha}{cucu}{78}
\DeclareMathOperator{\RR}{\mathbb R}%{\mathalpha}{cucu}{82}
\DeclareMathOperator{\QQ}{\mathbb Q}%{\mathalpha}{cucu}{81}
\DeclareMathOperator{\ZZ}{\mathbb Z}%{\mathalpha}{cucu}{90}
%%%%%%%%%%%%%%%%%%%%%%%%%%%%%%%%%%%%%%%%%%%%%%%%%%%%%%
\DeclareMathOperator{\sen}{sen}
\DeclareMathOperator{\arcsen}{arcsen}
\DeclareMathOperator{\senh}{senh} \DeclareMathOperator{\fr}{Frac}
%%%%%%%%%%%%%%%%%%%%%%%%%%%%%%%%%%%%%%%%%%%%%%%%%%%%
\newcommand{\ran}[1]{{\em Im\/}(#1)}
\newcommand{\ca}[1]{{\em card\/}(#1)}
\newcommand{\sop}[1]{{\em supp\/}(#1)}
\newcommand{\tra}[1]{{\em traza\/}(#1)}
\newcommand{\dis}[1]{d (#1)}
\newcommand{\inte}[1]{{\em Int\/}(#1)}
\newcommand{\conv}[1]{{\em conv\/}(#1)}
\newcommand{\too}{\longrightarrow}
\newcommand{\oo}{\overline}
\newcommand{\cq}{\begin{flushright} $\Box$ \end{flushright}}
\newcommand{\cqd}{\, \, \rule{7pt}{7pt}}
\def\limsup{\mathop{\overline{\rm lim}}}
\def\liminf{\mathop{\underline{\rm lim}}}
\renewcommand{\theenumii}{\arabic{enumii}}
\renewcommand{\labelenumii}{\theenumi.\theenumii.}
\renewcommand{\theenumiii}{\arabic{enumiii}}
\renewcommand{\labelenumiii}{\theenumi.\theenumii.\theenumiii}
%%%%%%%%%%%%%%%%%%%%%%%%%%%%%%%%%%%%%%%%%%%%%%%%%%%%%%%%%%%%%%%%%
\begin{document}
El estudio de los espacios de Sobolev con pesos es un tema de actualidad en el ámbito de la Teoría de Aproximación y polinomios ortogonales y con importantes aplicaciones numéricas. Este Trabajo Fin de Grado se enmarca en una investigación sobre  la localización de los ceros de los polinomios ortogonales de Sobolev con derivadas de primer orden.  A diferencia del caso estándar,
los ceros pueden salir de la envoltura convexa del soporte de la medida asociada. En el mismo se estudia  el comportamiento  del conjunto de ceros de los polinomios ortogonales de Sobolev  a partir de las matrices infinitas de momentos que definen productos escalares hilbertianos en los espacios de polinomios.

El objetivo principal  del trabajo es la utilización de algoritmos  para el análisis y  la visualización del conjunto de ceros de los polinomios para obtener resultados que arrojen luz sobre el problema de localización de estos conjuntos y  la relación de los soportes con la estructura del producto de Sobolev. Para ello se ha desarrollado un marco teórico unificado, extendiendo la teoría clásica, analizando y explorando  la
conexión de  la acotación de los ceros con  parámetros espectrales como la norma del operador de multiplicación y los valores singulares de ciertas matrices asociadas.\\ 

Durante el trabajo, el análisis teórico y la experimentación numérica han surgido algunas conjeturas y resultados  que se enmarcan en una  investigación en curso  que constituyen una de las aportaciones principales.\\


\end{document}